%-------------------------------%
%  Author: Alessandro Sciarra   %
%    Date: 11 Jul 2025          %
%-------------------------------%

%===================================================================================================%
\subsection{Test driven development as discipline}
%~~~~~~~~~~~~~~~~~~~~~~~~~~~~~~~~~~~~~~~~~~~~%
\AlertFrame{It is a discipline}
%~~~~~~~~~~~~~~~~~~~~~~~~~~~~~~~~~~~~~~~~~~~~%
\begin{frame}{It is a discipline}
    \begin{itemize}
        \item \alert{Do not bring it immediately into your daily work}
        \item Learn it in small exercise projects:
        \begin{itemize}
            \item Implement a \;\PS{\texttt{Stack}}\; object
            \item Implement a file parser
            \item Write a 1D integration tool with few different algorithms
            \item Simulate the 2D Ising model with a Monte Carlo algorithm
            \item Implement a simple game (e.g.\ hangman, minesweeper)
            \item{} [...]
        \end{itemize}
        \item \PB{Convince yourself about it}
        \item Get a feeling and choose ``your way''
    \end{itemize}
    \begin{varblock}{example}[0.78\textwidth]{Finally}
        \bfseries Get rid of the ``there-is-no-time-to-write-tests'' attitude!\\
    \end{varblock}
\end{frame}
%~~~~~~~~~~~~~~~~~~~~~~~~~~~~~~~~~~~~~~~~~~~~%
\begin{frame}{The purist original approach}
    \begin{varblock}{quote}[\textwidth]{Test driven development main rules}[Robert C. Martin]
        \begin{enumerate}[<+(1)->]
            \item You're not allowed to write any production code until you have first written a test that fails because the production code doesn't exist.
            \item You're not allowed to write more of a test than is sufficient to fail.
            And not compiling is failing.
            \item You're not allowed to write any more production code than is sufficient to pass the currently failing test.
        \end{enumerate}
    \end{varblock}
    \begin{varblock}{example}[0.6\textwidth]{}<5>
        \PS{How would it be a world in which everything always worked one minute ago?}
    \end{varblock}
\end{frame}
%~~~~~~~~~~~~~~~~~~~~~~~~~~~~~~~~~~~~~~~~~~~~%
\begin{frame}[fragile]{Writing the factorial in TDD}
    \vspace{-2mm}
    \begin{overlayarea}{\textwidth}{0.82\textheight}
        \begin{columns}[T]
            \begin{column}{0.49\textwidth}
                \begin{varblock}{example}[\textwidth]{Source file}<only@1-2>
                    \begin{lstlisting}[style=MyCpp]
                        // Empty file
                    \end{lstlisting}
                \end{varblock}
                \begin{varblock}{example}[\textwidth]{Source file}<only@3>
                    \begin{lstlisting}[style=MyCpp, belowskip=-9mm]
                        void factorial(int n){}
                    \end{lstlisting}
                    {\scriptsize\alert{Notice the function signature!}}
                \end{varblock}
                \begin{varblock}{example}[\textwidth]{Source file}<only@4-5>
                    \begin{lstlisting}[style=MyCpp, belowskip=-7mm]
                        void factorial(int n)
                        {
                          throw std::logical_error();
                        }
                    \end{lstlisting}
                    \only<4>{\scriptsize\alert{Still no need to return anything!}}
                \end{varblock}
                \begin{varblock}{example}[\textwidth]{Source file}<only@6-7>
                    \begin{lstlisting}[style=MyCpp]
                        int factorial(int n)
                        {
                          if(n<0)
                            throw std::logical_error();
                          else
                            return 1;
                        }
                    \end{lstlisting}
                \end{varblock}
                \begin{varblock}{example}[\textwidth]{Source file}<only@8-9>
                    \begin{lstlisting}[style=MyCpp]
                        int factorial(int n)
                        {
                          if(n<0)
                            throw std::logical_error();
                          else if(n==1)
                            return 1;
                          else
                            return n * factorial(n-1);
                        }
                    \end{lstlisting}
                \end{varblock}
                \begin{varblock}{example}[\textwidth]{Source file}<only@10-12>
                    \begin{lstlisting}[style=MyCpp]
                        int factorial(int n)
                        {
                          if(n<0)
                            throw std::logical_error();
                          else if(n<2)
                            return 1;
                          else
                            return n * factorial(n-1);
                        }
                    \end{lstlisting}
                \end{varblock}
                \begin{varblock}{example}[\textwidth]{Source file}<only@13>
                    \begin{lstlisting}[style=MyCpp]
                        #include |+<limits>+|

                        int factorial(int n)
                        {
                          if(n<0)
                            throw std::logical_error();
                          else if(n<2)
                            return 1;
                          else if(n>=13)
                            return INT_MAX;
                          else
                            return n * factorial(n-1);
                        }
                    \end{lstlisting}
                \end{varblock}
            \end{column}
            \hspace*{-4mm}
            \begin{column}{0.48\textwidth}
                \begin{varblock}{example}[\textwidth]{Test file}<only@1>
                    \begin{lstlisting}[style=MyCpp, belowskip=-9mm]
                        #include |+<vir/test.h>+|
                        #include |+"factorial.h"+|

                        TEST(Nothing){}
                    \end{lstlisting}
                    {\scriptsize\PS{It passes}}
                \end{varblock}
                \begin{varblock}{example}[\textwidth]{Test file}<only@2-4>
                    \begin{lstlisting}[style=MyCpp, belowskip=-10mm]
                        #include |+<vir/test.h>+|
                        #include |+"factorial.h"+|

                        TEST_CATCH(Negative_factorial,
                                   std::logical_error)
                        {
                          factorial(-17);
                        }
                    \end{lstlisting}
                    {
                        \scriptsize
                        \only<2>{\alert{It doesn't compile}}
                        \only<3>{\alert{It fails}}
                        \only<4>{\PS{It passes}}
                    }
                \end{varblock}
                \begin{varblock}{example}[\textwidth]{Test file}<only@5-6>
                    \begin{lstlisting}[style=MyCpp, belowskip=-10mm]
                        #include |+<vir/test.h>+|
                        #include |+"factorial.h"+|

                        TEST_CATCH(Negative_factorial,
                                   std::logical_error)
                        {
                          factorial(-17);
                        }
                        TEST(FactorialTest) {
                          COMPARE(1, factorial(1));
                        }
                    \end{lstlisting}
                    {\scriptsize\alt<5>{\alert{It fails}}{\PS{It passes}}}
                \end{varblock}
                \begin{varblock}{example}[\textwidth]{Test file}<only@7-8>
                    \begin{lstlisting}[style=MyCpp, belowskip=-10mm]
                        #include |+<vir/test.h>+|
                        #include |+"factorial.h"+|

                        TEST_CATCH(Negative_factorial,
                                   std::logical_error)
                        {
                          factorial(-17);
                        }
                        TEST(FactorialTest) {
                          COMPARE(1, factorial(1));
                          COMPARE(120, factorial(5));
                        }
                    \end{lstlisting}
                    {\scriptsize\alt<7>{\alert{It fails}}{\PS{It passes}}}
                \end{varblock}
                \begin{varblock}{example}[\textwidth]{Test file}<only@9-10>
                    \begin{lstlisting}[style=MyCpp, belowskip=-10mm]
                        #include |+<vir/test.h>+|
                        #include |+"factorial.h"+|

                        TEST_CATCH(Negative_factorial,
                                   std::logical_error)
                        {
                          factorial(-17);
                        }
                        TEST(FactorialTest) {
                          COMPARE(1, factorial(1));
                          COMPARE(120, factorial(5));
                          COMPARE(1, factorial(0));
                        }
                    \end{lstlisting}
                    {\scriptsize\alt<9>{\alert{It fails}}{\PS{It passes}}}
                \end{varblock}
                \begin{varblock}{example}[\textwidth]{Test file}<only@11>
                    \begin{lstlisting}[style=MyCpp, belowskip=-10mm]
                        #include |+<vir/test.h>+|
                        #include |+"factorial.h"+|

                        TEST_CATCH(Negative_factorial,
                                   std::logical_error)
                        {
                          factorial(-17);
                        }
                        TEST(FactorialTest) {
                          COMPARE(1, factorial(1));
                          COMPARE(120, factorial(5));
                          COMPARE(1, factorial(0));
                          COMPARE(479001600,
                                  factorial(12));
                        }
                    \end{lstlisting}
                    {\scriptsize\PS{It still passes}}
                \end{varblock}
                \begin{varblock}{example}[\textwidth]{Test file}<only@12>
                    \begin{lstlisting}[style=MyCpp, belowskip=-10mm]
                        #include |+<vir/test.h>+|
                        #include |+"factorial.h"+|
                        #include |+<limits>+|

                        TEST_CATCH(Negative_factorial,
                                   std::logical_error)
                        {
                          factorial(-17);
                        }
                        TEST(FactorialTest) {
                          COMPARE(1, factorial(1));
                          COMPARE(120, factorial(5));
                          COMPARE(1, factorial(0));
                          COMPARE(479001600,
                                  factorial(12));
                          COMPARE(INT_MAX,
                                  factorial(13));
                        }
                    \end{lstlisting}
                    {\scriptsize\alert{It fails}}
                \end{varblock}
                \begin{varblock}{example}[\textwidth]{Test file}<only@13>
                    \begin{lstlisting}[style=MyCpp, belowskip=-10mm]
                        #include |+<vir/test.h>+|
                        #include |+"factorial.h"+|

                        TEST_CATCH(Negative_factorial,
                                   std::logical_error)
                        {
                          factorial(-17);
                        }
                        TEST(FactorialTest) {
                          COMPARE(1, factorial(1));
                          COMPARE(120, factorial(5));
                          COMPARE(1, factorial(0));
                          COMPARE(479001600,
                              factorial(12));
                          COMPARE(INT_MAX,
                              factorial(13));
                        }
                    \end{lstlisting}
                    {\scriptsize\PS{It passes}}
                \end{varblock}
            \end{column}
        \end{columns}
        \FrameRemark<12->{In the code \;\CPP|INT_MAX|\; has been used for space reasons. The C++11 limit \;\CPP|std::numeric_limits<int>::max()|\; should be preferred, though.}
    \end{overlayarea}
\end{frame}
%~~~~~~~~~~~~~~~~~~~~~~~~~~~~~~~~~~~~~~~~~~~~%
\begin{frame}{How strict should you be with yourself?}
    \begin{center}
        \begin{tikzpicture}
            \node[label={[anchor=north, font=\Large, label distance=-5mm]90:{Find your way!}}] (fig)
            {
\includegraphics[width=0.7\textwidth]{Find-your-way}};
        \end{tikzpicture}
        \par \alert{The more edits you do at once, the more dangerous it is!}
    \end{center}
\end{frame}
%~~~~~~~~~~~~~~~~~~~~~~~~~~~~~~~~~~~~~~~~~~~~%
\begin{frame}{Disciplines feel slow, they are not slow}
    \begin{center}
        \includegraphics[height=0.8\textheight]{TDD-not-slow}
    \end{center}
\end{frame}
%~~~~~~~~~~~~~~~~~~~~~~~~~~~~~~~~~~~~~~~~~~~~%
\begin{frame}{The Red -- Green -- Refactor approach}
    \begin{onlyenv}<1>
        \begin{varblock}{quote}[0.85\textwidth]{What Kent Beck proposed}[Ian Cooper]
            \normalsize
            \begin{description}
                \item[\alert{\textbf{Red:}}] Write a little test that doesn't work, and perhaps doesn't even compile at first.
                \item[\PS{\textbf{Green:}}] Make the test work quickly, committing whatever sins necessary in the process.
                \item[\PB{\textbf{Refactor:}}] Eliminate all of the duplication created in merely getting the test to work.
            \end{description}
        \end{varblock}
        \vspace{-2mm}
        \begin{varblock}{quote}[0.92\textwidth]{About testing the implementation}[Ian Cooper]
            Sometimes it is useful to probe the implementation details because they are difficult to understand.
            You can do that, but \textbf{delete the tests afterwards!}\\[-3mm]~
        \end{varblock}
    \end{onlyenv}
    \begin{onlyenv}<2>
        \begin{center}
            \includegraphics[width=0.9\textwidth]{TDD-cycle}
        \end{center}
    \end{onlyenv}
\end{frame}
%===================================================================================================%


%===================================================================================================%
\subsection{BDD as alternative}
%~~~~~~~~~~~~~~~~~~~~~~~~~~~~~~~~~~~~~~~~~~~~%
\AlertFrame{
    \begin{tabular}{c}
        It is not \textbf{Bug} Driven Development!\\ \Emoji[0.5]{sweat-smile}
    \end{tabular}
}
%~~~~~~~~~~~~~~~~~~~~~~~~~~~~~~~~~~~~~~~~~~~~%
\begin{frame}{Behaviour driven development in its original idea}
    \begin{overlayarea}{\textwidth}{0.82\textheight}
        \begin{itemize}
            \item Dan North's \URL[PB]{https://dannorth.net/introducing-bdd/}{original idea} stemmed from TDD
            \begin{itemize}[<only@1>]
                \item Test method names should be sentences
                \item An expressive test name is helpful when a test fails
                \item Determine the next most important behaviour
                \item Requirements are behaviour, too
                \item BDD provides a ``ubiquitous language'' for analysis
            \end{itemize}
            \item<2-> Use cases become stories and each test implements an acceptance criterion for the story
        \end{itemize}
        \vspace{-2mm}
        \setlength{\labelsep}{\fontdimen2\font} % https://tex.stackexchange.com/a/266760/128737
        \begin{varblock}{example}[0.68\textwidth]{Example of use case as a story}<only@2>
            \textbf{Title:} Customer withdraws cash\\
            \begin{description}[~~so that]
                \setlength{\itemsep}{0pt}
                \item[\alert{As a}] customer,\\
                \item[\alert{I want}] to withdraw cash from an ATM,\\
                \item[\alert{so that}] I don't have to wait in line at the bank.
            \end{description}
        \end{varblock}
        \begin{varblock}{example}[0.8\textwidth]{Example of story acceptance criterion}<only@3>
            \textbf{Scenario:} Account is overdrawn past the overdraft limit\\
            \begin{description}
                \setlength{\itemsep}{0pt}
                \item[\alert{Given}] the account is overdrawn\\
                \item[\alert{And}] the card is valid\\
                \item[\alert{When}] the customer requests cash\\
                \item[\alert{Then}] ensure a rejection message is displayed\\
                \item[\alert{And}] ensure cash is not dispensed\\
                \item[\alert{And}] ensure the card is returned
            \end{description}
        \end{varblock}
    \end{overlayarea}
\end{frame}
%~~~~~~~~~~~~~~~~~~~~~~~~~~~~~~~~~~~~~~~~~~~~%
\begin{frame}{TDD and BDD}
    \begin{center}
        
\includegraphics[width=0.62\textwidth]{Different-POV}\\[3mm]
        After all they are disciplines to deal with the same idea!
    \end{center}
\end{frame}
%===================================================================================================%


%===================================================================================================%
\subsection{Take-home message}
%~~~~~~~~~~~~~~~~~~~~~~~~~~~~~~~~~~~~~~~~~~~~%
\begin{frame}{Wrapping up the advice}
    \vspace{-3mm}
    \begin{enumerate}[<+->]
        \item \textcolor{bg!80!normal text.fg}{«I do not have time to write tests»} \textbf{does not make any sense!}
        \item \PP{Writing tests is an inseparable part of the code development}
        \item \alert{Tests are the only way to change code still ensuring its correctness}
        \item \PS{Test the behaviour, not the implementation}
        \item Depending on the size of the project, you can go one or another way
        \item If you do not want to strictly do TDD, still envision your tests and API usage before implementing any code
        \item Now you are aware of how writing code should be done!
        \item Testing (and hence writing) code is doing science
    \end{enumerate}
\end{frame}
%~~~~~~~~~~~~~~~~~~~~~~~~~~~~~~~~~~~~~~~~~~~~%
\begin{frame}{And now?}
    \begin{uncoverenv}<2>
        \begin{center}
            
\includegraphics[width=0.68\textwidth]{Wow}\\[3mm]
            If this is not your reaction right now, it will be soon!
        \end{center}
    \end{uncoverenv}
\end{frame}
%===================================================================================================%